% =====================================================================
%  CH. 2 CF1 · CH. 3 CF2 · CH. 4 CF3 — LES TROIS GÉNÉRATIONS
% =====================================================================
\section{CF1 — le coffre à phrase (D75)}\label{sec:cf1}

\subsection{L'idée, et pourquoi elle a suffi un temps}

CF1 est né d'une demande du développeur : « stocker toutes les phrases maîtresses par
projet ». L'idée est la plus ancienne de la cryptographie symétrique : une phrase que
seul le maître connaît, dérivée en clé par une fonction coûteuse, protège tout le reste.
Elle a suffi un temps parce que la menace modélisée était simple — le \emph{vol au
repos} : un tiers qui s'empare du disque ou du dépôt ne doit rien pouvoir lire. Contre
cette menace, une phrase forte dérivée par Argon2id est une réponse complète.

Ce qui a fait grandir le coffre, c'est une autre menace, venue du terrain du
développeur : l'\emph{expropriation} (D84) — un tiers qui s'empare de la machine
\emph{pendant qu'elle vit}, phrase comprise. Contre elle, une phrase seule ne protège
pas : il faut que la clé n'existe nulle part entière. C'est l'histoire des deux
générations suivantes — CF1 n'est pas une erreur, c'est le socle : même CF3 dérive
K\_user par la même Argon2id aux mêmes planchers.

\subsection{La cérémonie de la clé}

\begin{lstlisting}
phrase (choisie par le maitre, D110 : alphanumerique, sans plafond)
  -> Argon2id(phrase, sel_argon 16 o, t=65536 Kio, m=3 it, p=1) = K_user 32 o
  -> HKDF-SHA256(K_user, sel_hkdf, info "coffre=v1") = KEK 32 o
  -> XChaCha20-Poly1305(JSON base, KEK, nonce 24 o frais, AAD = en-tete clair)
\end{lstlisting}

Le sel est propre au coffre (posé à l'init, 16 octets) ; le nonce est frais à
\emph{chaque écriture} ; l'empreinte du chiffretexte (sha256, posée dans l'en-tête)
nomme la corruption avant l'AEAD. La base claire est le JSON commun aux trois
générations :

\begin{lstlisting}
{ "format": ..., "protocole": ..., "cree": ..., "modifie": ...,
  "description": ..., "projets": { "<nom>": { "phrase": ..., ... }, ... } }
\end{lstlisting}

Chaque projet du dépôt y vit avec sa phrase maîtresse et ses métadonnées. Les gestes :
\code{init}, \code{lister}, \code{obtenir} (vers l'écran ou vers un fichier 0600),
\code{ajouter}, \code{retirer}, \code{verifier}, \code{changer-phrase},
\code{exporter}/\code{importer} (continuité), \code{bruler} (destruction manuelle).
L'autotest de falsification (19 contrôles à la naissance, D75) a établi la règle de la
maison : mauvaise phrase refusée ; mutation d'UN octet — chiffretexte, sel, nonce,
empreinte — chacune refusée nommément ; \code{lister} ne divulgue aucune phrase ;
l'export round-trip est exact.

\subsection{La garde de la phrase (D81) — le fichier est la remise}

L'épisode fondateur mérite d'être raconté, parce qu'il explique toute la discipline
actuelle. La rotation D80 avait produit une phrase neuve dont la remise — prévue « en
fin de tour » par affichage unique — n'a jamais eu lieu : la session d'exécution est
morte avant, d'une panne d'infrastructure de plus. La phrase ne vivait plus que dans le
bac à sable éphémère. Un décès d'infra de plus, et le coffre entier devenait
indéchiffrable à jamais. L'arbitrage est consigné : le risque de \textsc{perte} (vécu,
répété) dépasse le risque d'\textsc{exposition} (URL non devinable). Depuis : dès sa
production — init, rotation, changement — la phrase est écrite dans
\code{PHRASE-CONTROLE-COFFRE.txt} (une ligne, rien d'autre), commitée au dépôt, poussée
(GitHub + miroir Drive par le hook) et posée au racine du coffre Drive, \emph{avant
toute autre action de la session}. Une phrase entrée dans l'histoire git y demeure à
vie : on ne compte plus sur sa confidentialité, la rotation reste le remède.
La transition vers une garde plus dure est \emph{préparée} (D219) : le
séquestre chiffré D183 — la phrase répartie en cinq parts à seuil trois
(cf. § \ref{sec:cerd183}), la cérémonie au maître, la garde D81 demeure
ACTIVE jusqu'au geste vécu.

\section{CF2 — la clé répartie (D86)}\label{sec:cf2}

\subsection{L'idée : le secret qui n'existe nulle part}

CF2 répond à la question que CF1 laisse ouverte : que faire si l'adversaire tient la
machine \emph{vivante} ? La réponse de D86 : découper la clé maîtresse en fragments
tels qu'aucun supports ne suffise, et armer une sentinelle qui brûle les fragments si
un processus non autorisé lit la couche C. Le secret maître \textbf{M} (32 octets du
CSPRNG) n'est jamais écrit entier — nulle part, jamais : il vit reconstitué en RAM le
temps d'un déverrouillage, vérifié, puis effacé.

\subsection{Le découpage Shamir, en mots}

Un secret peut être \emph{réparti} : on le découpe en N fragments tels que n'importe
quels T d'entre eux le reconstituent, et que T$-1$ n'en disent \emph{rien} — pas
« difficile à calculer », \emph{rien} : la sécurité est informationnelle. Le mécanisme
(Shamir sur le champ GF($2^8$), polynôme générateur 0x11B — celui d'AES) : pour chacun
des 32 octets de M, un polynôme distinct de degré $T-1$ dont les coefficients sont
frais du CSPRNG ; le fragment $j$ porte les évaluations de ces polynômes au point
$x_j$. T fragments régressent les polynômes (Lagrange en $x=0$) et retrouvent M ;
T$-1$ fragments laissent tout M' possible — aucune discrimination, aucune attaque.

L'arbitrage terrain : \textbf{N = 5, T = 3}. Cinq gardes machine (répertoires distincts,
un fragment chacun, mode 0700) ; la relève exige \textbf{T copies hors-ligne} — trois
à seuil 3 : deux copies laisseraient le coffre mort si les gardes machine brûlaient
(l'autotest le prouve par restauration partielle refusée). Brûler les cinq gardes
machine = mort du coffre ; le code volé reste inexploitable (les enveloppes de carte
sont mortes avec lui) ; restaurer trois copies hors-ligne = résurrection. C'est un
contrat clair : \emph{ce que vous acceptez de perdre} contre \emph{ce que vous refusez
de voir exfiltré}.

\subsection{Les fragments — le format qui se transporte sur papier}

\begin{lstlisting}
TRIADE-CF2-FRAG v1
coffre: <sel HKDF du conteneur, 64 hex - l'identite STABLE du coffre>
x: <indice, decimal>
seuil: t=3/n=5
fragment: <64 hex - 32 o>
tag: <16 hex - sha256("TRIADE-CF2-frag" || x || fragment) tronque>
\end{lstlisting}

Le fragment est ASCII autodescriptif : il se lit, s'imprime, se recopie — l'exfiltration
d'un fragment est un vol physique, hors du modèle. Il est lié au \emph{sel} du coffre
(l'identité stable, posée à l'init) et jamais à l'empreinte de l'en-tête (le nonce
change à chaque écriture ; une liaison à l'en-tête invaliderait les copies papier à
chaque mutation — la faute a été dénoncée par le parcours CLI complet avant d'être
commise, D106). Le \code{tag} nomme le fragment fautif sans le divulguer : corruption,
mélange de fragments d'un autre coffre (sel divergent), collision d'indice — refus
nommés à l'import comme à la reconstitution.

\begin{terrain}[— les cinq originaux, la feuille et les copies (D131)]
La question du terrain, 31/10 : « la cérémonie m'a produit un fragments-cf2.txt et le
système attend 3 fragment-*.cf2 ». Les CINQ originaux vivent sous
\code{.keys\textbackslash gardes-cf2\textbackslash g1...g5\textbackslash
fragment-N.cf2f} — un fichier par garde, dans un dossier caché du clone. La FEUILLE
est l'export papier des TROIS fragments g3--g5 (les tags sont vérifiés avant
impression) : c'est une empreinte humaine, à ranger au sachet. Les COPIES USB sont les
fichiers \code{.cf2f} eux-mêmes, copiés hors-ligne : c'est eux que le contrôle de
relecture relit. La feuille et les copies ne sont pas le même objet — l'un se recopie
à la main, l'autre se restaure.
\end{terrain}

\subsection{Le conteneur, la sentinelle, la porte}

Le conteneur \code{TRIADE-CF2} dérive sa clé de M (HKDF, contexte \code{coffre=v2} —
nul chevauchement avec CF1 ni avec les fragments de carte) ; le champ \code{verif}
(16 octets) vérifie toute reconstitution avant que l'AEAD ne parle. La
\textbf{sentinelle} surveille la couche C (jamais les gardes ni le conteneur) à
\emph{double seuil} : un processus hors liste blanche \emph{et} une lecture de
contenu — deux signaux indépendants, parce qu'une fausse alerte serait une perte
auto-infligée. Moteur laboratoire : inotify + attribution \code{/proc} ; moteur de
production Windows : SACL + journal 4663 ou ETW (exigence consignée). La
\textbf{PORTE D120} — née de la demande du terrain : « si le déverrouillage est tenté
et qu'elle n'est pas lancée, tous nos efforts sont inutiles » — exige une sentinelle
vivante pour déverrouiller ; sur plateforme sans moteur, la porte \emph{avertit
nommément} à chaque déverrouillage. Le chapitre 5 raconte la fin de cette histoire.

La \textbf{liste blanche} de la sentinelle mérite son mot : elle blanchit par motifs
regex sur l'image de processus résolue \emph{et} sur la ligne de commande complète —
jamais « tout python3 » : c'est la ligne qui distingue \code{build\_m4.py} d'un
\code{-c} jetable. La liste est MINIMALE par construction (chaînes de build,
vérificateurs, l'éditeur du maître, la sentinelle elle-même) et testée par
falsifiabilité : un motif trop large blanchirait l'adversaire — c'est le défaut que
la suite prouve ne pas avoir.

La \textbf{destruction} (\code{bruler}) : trois passes d'écrasement (aléa, 0xFF, 0x00),
troncature, renommage aléatoire, unlink, fsync — et un \textbf{procès-verbal} nommant
le déclencheur (pid, image, fichier lu). Le conteneur et les gardes vides restent sur
place : le PV dit la mort, la vérification la confirme. La \textbf{restauration}
(\code{restaurer}) est atomique : chaque fragment est vérifié (tag + sel + indice
libre) et toute anomalie laisse \emph{rien} s'écrire — une remise partielle serait un
coffre menteur.

\subsection{Les critères de sortie de CF2 — et leurs preuves}

La spec scellée lie chaque promesse D86 à une preuve mesurable
(\code{v06/test\_cf2.py}) ; le maître peut les re-voir à volonté :

\begin{table}[htbp]
\centering\small
\caption{Critères de sortie D86 et leurs preuves (CF2).}
\label{tab:sortiecf2}
\begin{tabularx}{\textwidth}{@{}Y Y@{}}
\toprule
\textbf{Critère D86} & \textbf{Preuve} \\
\midrule
Shamir : reconstitution, seuil, refus nommés & T1 — KAT + 1000 vecteurs, interop
exhaustive C(5,3) \\
Sous T : informationnellement mort & T2 — un M' arbitraire accepté par les fragments
seuls : aucune discrimination \\
Conteneur : round-trip + fail-closed & T3 — en-tête, AEAD, \code{verif}, empreinte :
refus nommés \\
Exfiltration simulée $\to$ fragments détruits, indéchiffrable & F1 — sentinelle en
sous-processus, lecture non blanchie, PV + mort du coffre, restauration $\to$ vie : le
cycle complet \\
Fausses alertes NULLES en essai nominal & T4 — opérations blanches (build,
vérificateur simulés) : journal « blanchi », zéro brûlage \\
Latence détection $\to$ destruction bornée & F1 — mesurée, assertée sous la borne
consignée \\
Remise douteuse refusée sans écrire & F2 — atomicité, octet pour octet \\
Santé de la sentinelle & T5 — heart-beat frais, mort visible \\
\bottomrule
\end{tabularx}
\end{table}

\section{CF3 — la serrure piégée (D121)}\label{sec:cf3}

\subsection{L'idée : le témoin qui ne peut pas mentir}

CF3 est né de deux demandes du terrain, servies ensemble : « plutôt qu'une sentinelle
externe, une serrure du coffre piégée » et « une vulnérabilité réside durant le temps
où la session est ouverte ». L'idée : couper la clé en trois éléments dont l'un vit
\emph{hors machine} — le piège. Sans lui, le déverrouillage n'est pas refusé, il est
\emph{impossible} : il manque une moitié de la serrure. Et comme le piège est hors
machine, aucun état de machine falsifiable ne peut le remplacer : c'est un témoin que
le logiciel ne peut ni imiter ni truquer.

\subsection{Les trois éléments et l'époque}

\begin{itemize}
\item \textbf{K\_user} — la clé de l'utilisateur, \code{Argon2id(phrase,
sel\_argon)} : \emph{invariable à travers les époques} (la demande du terrain : « clé
utilisateur qui ne change pas »). Le sel et les paramètres sont posés à l'init et
jamais retouchés.
\item \textbf{S\_serrure} — la moitié « serrure » : 32 octets frais du CSPRNG à
\emph{chaque époque}, logée dans l'en-tête clair du conteneur (côté machine).
\item \textbf{S\_piège} — la moitié « piège » : 32 octets frais à chaque époque,
logée \emph{exclusivement sur le média séparé} — jamais écrite sur la machine hors
l'instant de l'acte.
\end{itemize}

L'\textbf{époque} (u64, posée à 1, incrémentée à chaque clôture) est l'horloge de la
métamorphose : chaque clôture renouvelle les moitiés, si bien qu'un secret volé ne
sert qu'une époque. \textbf{H\_projet} (sha256 de l'état scellé) lie la clé de clôture
au contenu — la course de clôture, § \ref{sec:course}.

\subsection{La dérivation et le champ qui ferme l'essai-hors-ligne}

\begin{lstlisting}
KEK = HKDF-SHA256(ikm = K_user || S_serrure || S_piege ;
                  salt = sel_hkdf ;
                  info = "triade-CF3/v1|" || epoque(le64) || H_projet)
verif = sha256(HKDF(ikm = K_user || V_piege ;
                    info = "triade-CF3/verif/v1"))[:16]
chiffretexte = XChaCha20-Poly1305(clair, KEK, nonce 24 o,
                                  AAD = en-tete jusqu'au nonce, 168 o)
\end{lstlisting}

Le champ \code{verif} est subtil et il mérite une phrase : le voleur de la
\emph{machine seule} ne peut même pas \emph{tester} des phrases contre le conteneur —
il lui manque V\_piège, la valeur publique qui vit sur le média. Le brute-force
hors-ligne n'est pas difficile, il est \emph{sans objet}. Contre le voleur de tout
(machine + média), la défense est Argon2id — borne consignée, même famille que CF1.

\subsection{Le fichier piège — un ASCII par époque}

\begin{lstlisting}
TRIADE-CF3-PIEGE v1
epoque: <u64 decimal>
projet: <H_projet hex (64)>
piege: <S_piege hex (64)>
verif: <V_piege hex (32)>
tag: <hex (16)>
\end{lstlisting}

Un fichier par époque, sur le média du sachet : \code{piege-CF3-epoque-1.txt},
\code{-2.txt}, \ldots{} Une clôture \textbf{ajoute} le piège N+1 et ne ronge jamais
les pièges passés — l'ancien piège est la relève du canonique précédent (si le média
perd le dernier piège, l'avant-dernier ne résurrect rien, mais il témoigne). Le
\code{tag} nomme la corruption du média sans divulguer ; un piège entièrement forgé se
présente comme une phrase erronée — bénin, refus nommé, jamais de destruction. Demain,
le média pourra être une clé USB à bouton physique (FIDO2) : seul le support change.

\subsection{Le cycle en deux actes — la fenêtre de session bornée}\label{sec:deuxactes}

\begin{enumerate}
\item \textbf{Acte I — ouverture.} Phrase + piège branché $\to$ sentinelle
instantanée (§ \ref{sec:sentinst}) $\to$ clair $\to$ \emph{brouillon} chiffré sur
disque sous une clé de session éphémère K\_session, qui ne vit qu'en RAM dans l'objet
Session. Puis K\_user, KEK et les moitiés sont effacés ; le média du piège est
\textbf{débranché} — geste humain, consigné au journal de cérémonie ; le conteneur
canonique n'est \emph{jamais} réécrit pendant la session.
\item \textbf{Session.} Lecture et écriture sur le brouillon (AEAD sous K\_session).
Une panne ou un abandon tue K\_session avec le processus : le brouillon devient un
déchet indéchiffrable, le canonique reste scellé. \emph{Une panne auto-nettoie.}
\item \textbf{Acte II — clôture.} Média rebranché, phrase \textbf{retapée} (une faute
de frappe est bénigne — aucun scellement n'a lieu). H\_final = sha256 du brouillon
final, checksum court de 8 hex \emph{affiché} : l'humain recoupe avec sa carte. Moitiés
fraîches, époque N+1, KEK de clôture dérivé \emph{avec} H\_final. Écriture du piège
N+1 (additive) puis du conteneur en \emph{staging} ; \textbf{re-lecture à froid} du
staging par la sentinelle complète — verte : remplacement atomique ; échec : staging
détruit, canonique et brouillon intacts, la clôture est rejouable.
\end{enumerate}

\subsection{La course de clôture et l'étiquette}\label{sec:course}

L'attaque : un malware attend la frappe de phrase à la clôture, vole la clé dérivée et
scelle \emph{son} état à la place. La parade est double. D'abord, la clé de clôture
est dérivée \emph{avec} H\_final : elle ne peut sceller que l'état exact dont le hash a
servi à la dériver. Ensuite, l'\textbf{étiquette} est vérifiée à l'ouverture : le clair
déchiffré doit hacher vers le H\_projet de l'en-tête. Le faux conteneur « étiquette
honnête, contenu substitué » est démasqué : sha256(contenu) $\neq$ H\_projet —
falsification \emph{prouvée}, destruction locale. Un KEK volé ne produit donc aucun
conteneur ouvrable pour un autre contenu, et le faux conteneur exige en outre d'écrire
le média du piège pendant qu'il est branché.

\subsection{Les trois témoins d'époque}\label{sec:temoins}

Le roll-back restaure un \emph{couple} cohérent (conteneur N + piège N) — l'AEAD
passerait seul. La défense est graduée : le \textbf{journal local}
(\code{.keys/journal-cf3.txt}) contrôle automatiquement à chaque acte (époque du
conteneur $\geq$ époque du miroir, sinon « RETOUR EN ARRIÈRE » nommé) mais un attaquant
admin peut l'effacer ; la \textbf{carte papier du sachet} consigne le checksum de 8 hex
à la main à chaque clôture — indécevable par logiciel ; le \textbf{journal Git}
(consigné à chaque cérémonie de clôture, poussé) est le témoin archivé hors machine.
Le verdict à froid (\code{verifier}) affiche le checksum pour le recoupement humain :
\emph{la machine dit, la carte confirme, vous seul jugez}.

\begin{terrain}[— ce que dit un verdict réel]
Extrait du rapport de migration du coffre réel (30/10/2026, boîte noire du GUI) — le
format que vous verrez à chaque \code{verifier} :
\begin{lstlisting}
[CF3] VERDICT : VERT - epoque 1, checksum 0cc89a96
    (carte papier / journal Git : recoupez), 1 projet(s)
\end{lstlisting}
Époque, checksum, projet(s) : les trois choses à recouper. Le checksum se note sur la
carte du sachet, l'époque se suit d'une clôture à l'autre, le compte de projets doit
être le vôtre.
\end{terrain}

\subsection{Les critères de sortie de CF3 — et leurs preuves}

Même loi (\code{v06/test\_coffre\_cf3.py}, la suite de falsifiabilité) :

\begin{table}[htbp]
\centering\small
\caption{Critères de sortie CF3 et leurs preuves.}
\label{tab:sortiecf3}
\begin{tabularx}{\textwidth}{@{}Y Y@{}}
\toprule
\textbf{Critère} & \textbf{Preuve attendue} \\
\midrule
Round-trip deux actes & init $\to$ ouvrir $\to$ modifier $\to$ fermer $\to$ ouvrir :
contenu exact, époque incrémentée, moitiés remplacées, \code{sel\_argon} invariant \\
Phrase erronée (acte I ET II) & refus nommé, JAMAIS de destruction, aucune écriture \\
Piège absent & déverrouillage impossible, refus nommé \\
Piège périmé / roll-back du couple & refus nommé + journal \\
Tag du piège invalide & refus nommé SANS destruction (média hors limite) \\
Falsification prouvée côté machine & destruction locale + PV nommé, média et fragments
intacts \\
Course de clôture & KEK volé + faux conteneur (étiquette menteuse OU honnête à
contenu substitué) : refus + destruction, jamais une ouverture \\
Panne simulée à chaque phase de clôture & canonique intact, cycle rejouable \\
Auto-nettoyage par panne & brouillon indéchiffrable sans K\_session, canonique jamais
réécrit pendant la session \\
Non-régression & sceau\_I verifier CODE 0, suites CF2/sceau-I/GUI vertes \\
\bottomrule
\end{tabularx}
\end{table}
